\documentclass{report}%
\usepackage[T1]{fontenc}%
\usepackage[utf8]{inputenc}%
\usepackage{lmodern}%
\usepackage{textcomp}%
\usepackage{lastpage}%
\usepackage{geometry}%
\geometry{a4paper=True,margin=25mm}%
\usepackage{amsmath}%
\usepackage{amssymb}%
\usepackage{amsfonts}%
\usepackage{mathtools}%
\usepackage{bm}%
\usepackage{physics}%
\usepackage{graphicx}%
\usepackage{microtype}%
\usepackage{iftex}%
\usepackage{listings}%
\usepackage{listingsutf8}%
\usepackage{xcolor}%
\usepackage[strings]{underscore}%
\usepackage{url}%
\usepackage{xurl}%
\usepackage{hyperref}%
\usepackage{bookmark}%
\usepackage[square,numbers]{natbib}%
%
\ifLuaTeX\usepackage{fontspec}\usepackage{unicode-math}\else\usepackage[utf8]{inputenc}\usepackage[T1]{fontenc}\fi\usepackage[czech]{babel}%
\graphicspath{{figures/}}%
\renewcommand{\bibname}{Použitá literatura}%
\renewcommand{\bibsection}{\chapter*{\bibname}\addcontentsline{toc}{chapter}{\bibname}}%
% Robust LaTeX settings for AutoGenBook (LuaLaTeX-friendly).
\ProvidesFile{robustness.tex}[2026/01/01 Robust LaTeX settings]

\RequirePackage{xparse}
\RequirePackage{fvextra}
\RequirePackage{enumitem}
\RequirePackage{adjustbox}
\RequirePackage{tabularx}
\RequirePackage{ltablex}
\keepXColumns

% Fallback for unknown environments from conversions.
\ProvideDocumentEnvironment{para}{+b}{\par\smallskip\noindent#1\par\smallskip}{}

% Prompt/code block (breaks anywhere to avoid overfull boxes).
% Prompt/code block with safe line breaking.
\DefineVerbatimEnvironment{PromptBlock}{Verbatim}{
  formatcom=\ttfamily\small,
  breaklines=true,
  breakanywhere=true,
  breaksymbolleft=,
  breaksymbolright=
}

% Make plain verbatim blocks breakable as well.
\DefineVerbatimEnvironment{verbatim}{Verbatim}{
  formatcom=\ttfamily\small,
  breaklines=true,
  breakanywhere=true,
  breaksymbolleft=,
  breaksymbolright=
}

% Listings defaults (robust line breaking).
\lstset{
  basicstyle=\ttfamily\small,
  breaklines=true,
  breakatwhitespace=false,
  columns=fullflexible,
  keepspaces=true,
  showstringspaces=false
}

% Description list labels should wrap to next line.
\setlist[description]{style=nextline,leftmargin=0pt,labelsep=0.5em}

% Safer line breaking for prose.
\emergencystretch=3em
\tolerance=4000
\pretolerance=1000
\hfuzz=0.2pt

% Better URL breaks.
\Urlmuskip=0mu plus 2mu

% Images never exceed linewidth.
\makeatletter
\def\maxwidth{%
  \ifdim\Gin@nat@width>\linewidth\linewidth\else\Gin@nat@width\fi}
\setkeys{Gin}{width=\maxwidth,keepaspectratio}
\makeatother

% Fit-width wrapper for wide tables.
\NewDocumentEnvironment{fitwidth}{+b}{\begin{adjustbox}{max width=\linewidth}#1\end{adjustbox}}{}

% LuaLaTeX fonts and Unicode fallbacks.
\ifLuaTeX
  \defaultfontfeatures{Ligatures=TeX}
  \newcommand{\AutogenMainFont}{Latin Modern Roman}
  \IfFontExistsTF{Latin Modern Roman}{}{\renewcommand{\AutogenMainFont}{TeX Gyre Termes}}
  \newcommand{\AutogenSansFont}{Latin Modern Sans}
  \IfFontExistsTF{Latin Modern Sans}{}{\renewcommand{\AutogenSansFont}{TeX Gyre Heros}}
  \newcommand{\AutogenMonoFont}{DejaVu Sans Mono}
  \newcommand{\AutogenMonoBold}{DejaVu Sans Mono Bold}
  \newcommand{\AutogenMonoItalic}{DejaVu Sans Mono Oblique}
  \newcommand{\AutogenMonoBoldItalic}{DejaVu Sans Mono Bold Oblique}
  \IfFontExistsTF{DejaVu Sans Mono}{}{
    \renewcommand{\AutogenMonoFont}{Latin Modern Mono}
    \renewcommand{\AutogenMonoBold}{Latin Modern Mono Bold}
    \renewcommand{\AutogenMonoItalic}{Latin Modern Mono Slanted}
    \renewcommand{\AutogenMonoBoldItalic}{Latin Modern Mono Bold Slanted}
  }

  \setmainfont{\AutogenMainFont}
  \setsansfont{\AutogenSansFont}
  \IfFontExistsTF{\AutogenMonoBold}{
    \setmonofont{\AutogenMonoFont}[
      BoldFont=\AutogenMonoBold,
      ItalicFont=\AutogenMonoItalic,
      BoldItalicFont=\AutogenMonoBoldItalic,
      Scale=MatchLowercase
    ]
  }{
    \setmonofont{\AutogenMonoFont}[Scale=MatchLowercase]
  }
\fi
%
\title{Mini průvodce AI ve výuce}%
\date{\today}%
%
\begin{document}%
\normalsize%
\lstset{backgroundcolor=\color[gray]{0.95},basicstyle=\ttfamily\small,breaklines=true,breakatwhitespace=false,columns=fullflexible,keepspaces=true,frame=single,numbers=left,numberstyle=\tiny,tabsize=4,inputencoding=utf8,extendedchars=true,showstringspaces=false}%
\hypersetup{pageanchor=false}%
\pagenumbering{roman}%
\maketitle%
\tableofcontents%
\clearpage%
\hypersetup{pageanchor=true}%
\pagenumbering{arabic}%
\chapter{Úvod}%
Přehled cílů průvodce zaměřený na kontext vysokého školství. Stručně vysvětlí základní pojmy AI a rozdíl mezi nástroji pro učitele (asistenti, autogradery, nástroje pro tvorbu obsahu) a nástroji pro studenty (adaptivní systémy, personalizované cesty). Nastíní hlavní otázky: jak AI zlepšuje učení a zpětnou vazbu, co je třeba zvážit při zavádění (pedagogické cíle, kvalita dat, ochrana soukromí), a jak hodnotit přínos. Úvod zmíní konkrétní, ověřitelné příklady z praxe (např. chatboty používané jako kursovní asistenti, autogradery v kurzech programování, adaptivní moduly v blended learningu) a připraví čtenáře na praktické kapitoly.%


%
Tento krátký průvodce je určen učitelům a administrátorům vysokých škol, kteří chtějí rychle a prakticky začít využívat nástroje umělé inteligence ve výuce; popisuje scénáře, ukázky z praxe, pedagogické přínosy a kroky pro bezpečné a etické zavedení do kurzů. (General background knowledge)
Stručná orientace v pojmech a typech nástrojů:
\begin{itemize}
  \item Nástroje pro učitele: asistenti pro přípravu obsahu, autogradery, generátory cvičení a multimédií. (General background knowledge)
  \item Nástroje pro studenty: adaptivní moduly, konverzační chatboty a systémy průběžné zpětné vazby. (General background knowledge)
\end{itemize}
Hlavní otázky průvodce:
\begin{itemize}
  \item Jak AI podporuje učení a zpětnou vazbu; jak definovat výukové cíle a metriky úspěchu? (General background knowledge)
  \item Jaké jsou požadavky na kvalitu a ochranu dat, souhlas studentů, inkluze a jak prakticky vyhodnocovat dopad? (General background knowledge)
\end{itemize}
Ilustrační příklady z praxe:
\begin{itemize}
  \item Chatboti jako kurzoví asistenti; autogradery pro hodnocení kódu; adaptivní moduly měnící obtížnost podle výkonu. (General background knowledge)
\end{itemize}
Praktické kapitoly obsahují scénáře nasazení s postupem a integrací do LMS, ověřené příklady, tipy pro piloty a škálování, a doporučení pro měření účinnosti, řízení rizik, ochranu údajů a etiku. (General background knowledge)%

%
\chapter{Praktické příklady}%
Konkrétní scénáře použití AI ve vysokoškolské výuce s popisem účelu, průběhu použití, požadovaných nástrojů, ověřených implementací na univerzitách, očekávaných přínosů a možných omezení. Každý scénář obsahuje krátký postup zavedení, potřebné zdroje a metody pro měření efektivity.
Hlavní podtémata (příklady fungujících řešení):
- Automatizovaná zpětná vazba a autogradery: použití pro programování (autograde systémy), testy s otevřenými odpověďmi (automatická kontrola struktur), rychlá formativní zpětná vazba, úspora času učitelům.
- Adaptivní a personalizované učení: moduly, které přizpůsobují obsah a úlohy podle výkonu studenta (příklady kurzů, kde adaptivní cesty zvýšily udržení znalostí), doporučení pro nastavení parametrů a integraci se LMS.
- Chatboti a virtuální asistenti pro studenty a učitele: praktické nasazení jako FAQ a 24/7 podpora, asistenti pro plánování lekcí a administrativu (příklad univerzitních nasazení chatbotů jako podpora učitelů), doporučení pro design dialogu a dohled člověka.
- Podpora hodnocení a korektury: nástroje pro rychlé zpracování a kategorizaci studentských odpovědí, pomoc při hodnocení esejí (doplňkové nástroje, které zrychlují práci hodnotitelů, ne nahrazují odborné posouzení).
- Learning analytics a podpora zásahů: dashboardy ukazující rizikové studenty, přehled angažovanosti v online aktivitách, jak využít data pro cílené intervence.
- Tvorba a adaptace učebních materiálů: nástroje pro generování cvičení, vytvoření studijních opor a multimedia obsahu při zachování akademické kvality.
Každé téma obsahuje krátkou ukázku implementace, konkrétní přínosy (time‑saving, lepší personalizace, vyšší angažovanost) a seznam rizik (bias, kvalita dat, přehnané spoléhání se na automatizaci).%
\section{Automatizovaná zpětná vazba a podpora hodnocení}%
Popis praktických nasazení autograderů a nástrojů pro rychlé zpracování studentských odpovědí. Obsahuje účel použití, krok‑za‑krokem postup zavedení pro programovací úlohy i otevřené odpovědi, požadované technologie a integrační body do LMS, ukázku krátké implementace, metody měření úspory času a kvality hodnocení a seznam běžných rizik a mitigací (např. chyby v testovacích sadách, potřeba lidského dohledu).%


%
(General background knowledge) Tato část popisuje praktické nasazení autograderů a kroky pro zavedení v kurzu.
(General background knowledge) Účel nasazení: zrychlit rutinní hodnocení, poskytovat okamžitou zpětnou vazbu a uvolnit čas vyučujícímu.
(General background knowledge) Základní kroky zavedení (krátký postup):
\begin{enumerate}
\item (General background knowledge) Definujte vhodné typy úloh a připravte testovací sady včetně okrajových případů; dělejte peer review testů.
\item (General background knowledge) Vyberte nástroj kompatibilní s LMS, integrujte ho do workflow odevzdání a pilotujte na malé skupině, sbírejte chyby a feedback.
\item (General background knowledge) Nastavte pravidlo lidského přezkumu pro sporné či otevřené odpovědi a iterujte podle zjištěných chyb.
\end{enumerate}
(General background knowledge) Měřte přínos (čas ručního hodnocení, opravy testů, spokojenost studentů, A/B piloty); mitigujte rizika (peer review testů, povinný lidský přezkum, testovací prostředí a IT podpora) a začínejte malými piloty.%

%
\section{Adaptivní a personalizované učení}%
Praktický průvodce pro nasazení adaptivních modulů: definování výukových cest, nastavení parametrů přizpůsobení výkonu studenta, požadavky na data a integraci se stávajícím LMS. Obsahuje krátký implementační scénář, doporučené kroky pro pilotní test, ukazatele efektivity (např. metriky učení a udržení), a popis možných omezení a zásad návrhu obsahu pro personalizaci.%


%
Úvodní poznámka: Níže shrnuté kroky a doporučení pro nasazení adaptivních modulů — praktické vodítko, které je třeba ověřit na místních datech a infrastruktuře (general background knowledge).
Cíl: nabídnout studentům individuální výukové cesty podle výkonu a potřeb, zvýšit efektivitu učení a udržení znalostí.
\begin{enumerate}
\item Definujte výukové cíle a rozdělte obsah do krátkých modulů s měřitelným mastery; \item Navrhněte větve (nápravné, rozšiřující, opakovací) a pravidla adaptace (prahy, opakování, spacing); \item Zajistěte datové zdroje, anonymizaci, integraci s LMS a pilotujte na malé skupině.
\end{enumerate}
Pilot: malá+kontrolní skupina, průběžné monitorování, rychlé iterace a dokumentace rozhodnutí pro škálování.
Měřítka: pre/post skóre, udržení, engagement (čas/aktivita), míra dokončení, podpora a kvalita reakcí.
Rizika a zásady: cold start a záložní cesta, riziko zkreslení a nutnost auditu, ochrana dat; návrh obsahu: krátké kroky, jasná kritéria mastery, vysvětlení doporučení a možnost přepnutí; šablona modulu: cíl, mapování, pravidla, data, pilot.%

%
\section{Chatboti a virtuální asistenti pro studenty a učitele}%
Návod pro praktické použití chatbotů jako FAQ, 24/7 podpory a asistentů pro administrativní úkony: návrh dialogů, definice rolí a pravidel eskalace k člověku, technické požadavky a způsoby nasazení v akademickém prostředí. Obsahuje příklad pracovního postupu pro zavedení, metody monitorování kvality odpovědí a doporučení pro dohled a ochranu soukromí.%


%
\begin{flushleft}
Účel: chatboty v akademickém prostředí slouží jako 24/7 FAQ, první linie podpory a asistent pro opakující se úkony; pedagogické přínosy při správném návrhu a monitorování.
Doporučený postup: 1) vymezit role a eskalace (citlivé údaje, změny známek => člověk), 2) sesbírat 20–50 scénářů a šablony dialogů, 3) pilot v LMS s fallback, 4) monitorovat logy, upravovat intenty a škálovat governance.
Návrh dialogů: používejte intent‑slot šablony, zobrazujte zdroj a možnost kontaktu; definujte prahy nejistoty pro okamžitou eskalaci.
Technické požadavky: integrace s LMS, autentizace a řízení přístupu, auditní logy, anonymizace a jasné zásady retence.
Monitorování: sledujte přesnost (lidský audit), míru eskalací, spokojenost uživatelů a zavádějte pravidelné lidské kontroly.
Etika a soukromí: transparentně informujte uživatele, nabídněte opt‑out a bezpečný kanál pro citlivé záležitosti.
Kontrolní seznam pro pilot: ověřené FAQ/odpovědi, pravidla eskalace a kontaktní osoby, ochrana dat a mechanismus auditu.
\end{flushleft}%

%
\section{Learning analytics a podpora cílených intervencí}%
Konkrétní scénář vytvoření dashboardu a pracovních postupů pro identifikaci rizikových studentů a plánování zásahů. Popis požadovaných datových zdrojů, návrh klíčových metrik angažovanosti, praktický postup nasazení analytiky a doporučené metody vyhodnocení zásahů. Součástí jsou také upozornění na etické a právní aspekty práce s daty.%


%
\noindent Poznámka: následující stručný scénář nasazení learning analytics a zásahů je general background knowledge a slouží jako šablona přizpůsobitelná místním podmínkám.
\noindent Účel: rychle identifikovat studenty s klesající angažovaností a spustit předem definované podpůrné zásahy (e‑mail, konzultace, doporučené materiály); pilot 4–8 týdnů, upravovat prahy.
\noindent Požadované datové zdroje: LMS (přihlášení, návštěvy, dokončení úkolů, diskuse), hodnocení, aktivita u e‑learningových objektů, administrativní údaje (jen eticky/legálně).
\noindent Doporučené metriky: frekvence přihlášení týdně, \% dokončených povinných aktivit, trend průběžného skóre, interakce s materiály.
\noindent Dashboard a upozornění: seznam rizikových studentů s filtry (kurz/skupina/období), krátký komentář, upozornění instruktorovi — ne plně automatické zásahy.
\noindent Plán zásahů: automatický e‑mail s nabídkou podpory, manuální ověření instruktorem, personalizovaná doporučení nebo schůzka s tutorem.
\noindent Vyhodnocení: měřte změny metrik a průběžných známek před/po, používejte A/B piloty nebo historické kontroly a krátké průzkumy spokojenosti.
\noindent Etika a checklist: sbírejte jen nezbytná data, informujte studenty a umožněte opt‑out, zajistěte bezpečné ukládání, omezení přístupu, retenci, schválení etiky a připravené šablony zásahů a eskalace.%

%
\section{Tvorba a adaptace učebních materiálů pomocí AI}%
Praktické postupy pro generování cvičení, tvorbu studijních opor a multimediálního obsahu s kontrolou akademické kvality: výběr nástrojů, workflow pro editaci a validaci vytvořeného materiálu, ukázková implementace pro jeden modul kurzu a metody měření přínosu (čas tvorby, zpětná vazba studentů). Upozornění na rizika spojená s kvalitou a autorskými právy a doporučení pro kontrolu obsahu.%


%
Účel a zásady
\begin{itemize}\item Cílem zrychlit tvorbu cvičení, studijních opor a multimédií při zachování akademické kvality a souladu s výukovými cíli; pracujte iterativně: AI navrhne, učitel upraví a validuje, pilot a úpravy.\end{itemize}
Výběr nástrojů — krátký checklist
\begin{itemize}\item Podpora vašich formátů (LMS, QTI/CSV, multimédia), transparentnost modelu a editovatelnost, řízení autorských práv/provenance a datové zabezpečení/nasazení on‑premise.\end{itemize}
Workflow (krok‑za‑krok)
\begin{enumerate}\item Definujte výukové cíle a formát; vytvořte šablonu promptu; generujte první verzi; učitel kontroluje správnost, shodu s cíli, obtížnost, inkluzi a rizika; pilot, sběr metrik a finalizace/integace do LMS.\end{enumerate}
Ukázka a kontrola kvality
\begin{itemize}\item Příklad: učitel nastaví cíle a šablony, AI vygeneruje varianty, náhodná kontrola; zajistit faktickou správnost, respekt k autorským právům, inkluzi, archivaci promptů a měření přínosu (čas, kvalita, spokojenost, výsledky).\end{itemize}%

%
\chapter{Zavedení do praxe, etika a hodnocení efektivity}%
Praktické kroky pro implementaci AI v kurzech: plánování (cíle výuky, metriky úspěchu), pilotní nasazení, školení učitelů, zapojení studentů a iterativní zlepšování. Doporučení pro tvorbu interních pravidel (governance), zajištění ochrany osobních údajů a transparentnosti (co nástroj dělá a jak jsou použita data). Sekce obsahuje rozhodovací strom pro výběr nástroje (úroveň technické podpory, soulad s LMS, dostupnost dat), návrhy metrik pro hodnocení přínosu (úspora času, zlepšení výsledků učení, spokojenost studentů) a příklady měření z praxe. Zvláštní pozornost věnována inkluzi (jak AI pomáhá adaptovat výuku pro různé potřeby) a rizikům (hnětení biasu, bezpečnost, akademická poctivost). Obsahuje také doporučení pro profesní rozvoj učitelů a zapojení fakultní administrativy.%


%
Praktické kroky pro zavedení AI do kurzu: jasné výukové cíle, definice automatizovaných částí a kontrolní body pro pilotní fázi \cite{node_2-5}.
Doporučený workflow: \begin{enumerate}
\item Definice cíle/formátu AI (jaké výstupy) \cite{node_2-5}.
\item Archivace promptů/konfigurací a lokální kontrola učitelem \cite{node_2-5}.
\item Pilotní nasazení s měřením (čas, kvalita, spokojenost, výsledky) a iterativní úpravy \cite{node_2-5}.
\end{enumerate}
Volba nástroje a metriky: podporované formáty (LMS, QTI/CSV, multimédia), transparentnost/editovatelnost, bezpečnost/on‑premise; metriky — úspora času, zlepšení výsledků, spokojenost, kvalita materiálů \cite{node_2-5}.
Etika a governance: transparentní komunikace o využití dat, uchovávání verzí promptů pro audit, interní pravidla a profesní rozvoj učitelů; monitorovat rizika (zkreslení, bezpečnost, akademická poctivost) \cite{node_2-5}.%

%
\bibliographystyle{unsrt}%
\bibliography{refs}%
\end{document}